\newpage
\section{Klasyczne modele sieci neuronowych}

Za początek konceptu sztucznych sieci neuronowych można przyjąć 1943 rok, w~którym to McCulloch i~Pitts, opierając się na działaniu komórki mózgowej (neuronu), opisali model matematyczny, który stał się podwalinami dzisiejszej sztucznej inteligencji. Komórka ta została przybliżona do matematycznego modelu będącego bramką logiczną. Neuron integruje docierające do niego sygnały w~swoim wnętrzu i~generuje dalszy sygnał, jeżeli energia przychodzącego sygnału przekroczy graniczną wartość \cite{raschka_python}.
Sztuczny neuron jest elementem procesującym informacje \cite{haykin2009neural} posiadający następujące cechy:

\begin{itemize}
    \item Neuron odbiera sygnały (informacje), tak jak jego biologiczny odpowiednik posiada synapsy, które są charakteryzowane przez wagi skalujące odbierany sygnał. W~przypadku sztucznych neuronów wagi te mogą przyjmować wartości ujemne.
    \item Neuron przetwarza odebrany sygnał wykorzystując funkcję sumującą, która dokonuje liniowej kombinacji wag i~sygnałów.
    \item Neuron wykorzystuje funkcję aktywacji do wyznaczenia ostatecznej wartości przekazywanego dalej sygnału. Sygnał również jest skalowany do odpowiednich zakresów.
\end{itemize} 
    
\textbf{Wielowarstwowe sieci neuronowe} (ang. Multi-Layer Perceptron – MLP) to zbiór sztucznych neuronów uporządkowanych w~warstwach. Każda warstwa przetwarza dane i~przekazuje je do kolejnej warstwy. Pierwsza warstwa to warstwa wejściowa, która przyjmue dane wejściowe i~przekazje sygnał, ostatnia to warstwa wyjściowa, która może mieć więcej niż jedno wejście i~wyjście, a~pomiędzy nimi znajdują się warstwy ukryte \cite{tutorial_nn_activation}.

\textbf{Perceptron Rosenblatta} jest szczególnym typem sztucznego neuronu, powszechnie nazywanym perceptronem. Perceptron jest najprostszą wersją sieci neuronowej, której zadaniem jest klasyfikacja wzorców będących liniowo rozdzielnymi, tj. wzorce, które leżą po przeciwnych stronach hiperpłaszczyzny. 
Jest to pojedynczy neuron wraz z~wagami na wejściu, w~tym z~wagą zerową (wyrazem wolnym). 
Aby perceptron poprawnie klasyfikował, klasy muszą być liniowo rozdzielne. W~przypadku zagadnienia dwuwymiarowego oznacza dwa obszary, które można rozdzielić prostą.
Ograniczenie to dyskwalifikuje użycie perceptronu Rosenblatta do analizy danych prognostycznych silników lotniczych ze względu na ich nieliniową złożoność, spowodowaną nieliniowymi procesami termodynamicznymi i~aerodynamicznymi, co można zauważyć na rysunku \ref{fig:scatter}.


\textbf{Funkcja aktywacji} definiuje wartość sygnału wyjścia neuronu w~funkcji liniowej kombinacji wejść i~wag. Jest matematycznym modelem odwzorowującym sygnał wejściowy na sygnał wyjściowy neuronu. W~przypadku perceptronu funkcją aktywacji jest funkcja skoku jednostkowego (ang. \textit{step function}) lub funkcja signum. Funkcje aktywacji można podzielić na dwie kategorie: liniowe i~nieliniowe. Funkcje liniowe są rzadko stosowane w~praktyce, ponieważ sieci neuronowe z~liniowymi funkcjami aktywacji mają ograniczoną zdolność do modelowania złożonych zależności. Nieliniowe funkcje aktywacji, takie jak sigmoidalna, ReLU (ang. \textit{Rectified Linear Unit}) czy tanh (tangens hiperboliczny), pozwalają sieciom neuronowym na modelowanie skomplikowanych wzorców i~relacji w~danych. Wybór odpowiedniej funkcji aktywacji ma kluczowe znaczenie dla efektywności uczenia się i~ogólnej wydajności sieci neuronowej.

Nieliniowość funkcji aktywacji sprawia, że sieć neuronowa przekształca się w~nieliniowe narzędzie aproksymacyjne. W~przypadku użycia liniowej funkcji aktywacji model sieci: 
\[f(x)= \beta_0 +\sum_{k=1}^{K}{\beta_k g(w_{k0}+\sum_{j=1}^{p}{w_{kj}X_j})}\],
gdzie: 
\begin{itemize}
    \item $g(\cdot)$ - funkcja aktywacji,
    \item $X_j$ - wejścia,
    \item $w_{kj}$ - wagi,
    \item $\beta_k$ - współczynniki liniowe,
    \item $K$ - liczba neuronów w~warstwie ukrytej,
    \item $p$ - liczba wejść,
\end{itemize}
redukuje się do postaci liniowej względem wejść $X_j$ \cite{islp}.

Przykładowe funkcje aktywacji:
\begin{description}
    \item[Krokowa (Heaviside)] Neuron wykorzystujący tę funkcję nazywany jest neuronem McCullocha-Pittsa. Jest to najprostsza funkcja aktywacji, która modeluje podstawowe zachowanie biologicznego neuronu. Jednakże jej zastosowanie w~praktyce jest ograniczone ze względu na brak różniczkowalności w~punkcie progowym, co utrudnia stosowanie metod gradientowych do trenowania sieci neuronowych \cite{haykin2009neural}. 
    \[Th_{\geq 0}(a) = \begin{cases}1, & a>0, \\0, & a~\leq 0\end{cases}\]
    \item[tanh (tangens hiperboliczny)] Funkcja ta jest podobna do sigmoidalnej, ale jej wartości wyjściowe mieszczą się w~zakresie od -1 do 1. Funkcja tanh jest różniczkowalna i~ma korzystne właściwości gradientowe, co czyni ją popularnym wyborem w~sieciach neuronowych. Wykorzystano ją w zaprojektowanych w niniejszej pracy autoenkoderach klasycznych jako funkcję warstw wyjściowych, aby przeskalować sygnały wyjściowe do zakresu [-1, 1], co ułatwia interpretację zrekonstruowanych parametrów fizycznych. 
    \[\tanh(a) = \frac{1-e^{-2a}}{1+e^{-2a}}\]
    \item[ReLU (Rectified Linear Unit)] Ułatwia tzw. rzadkie kodowanie (ang. \textit{sparse coding}), ze względu na niską liczbę aktywowanych neuronów w~sieci. Jej implementacja może prowadzić do zjawiska „martwych neuronów", gdzie wagi przestają się aktualizować, jeśli kombinacja liniowa wag i~wejść jest zawsze ujemna. Nie jest różniczkowalna w~punkcie zero, co może utrudniać trenowanie sieci neuronowej za pomocą metod gradientowych. 
    \[ReLU(a) = \max(0, a)\]
    Wykorzystano tę funkcję w~warstwach ukrytych zaprojektowanych w niniejszej pracy klasycznych sieci ze względu na efektywność obliczeniową.
    \item [Softplus] Gładkie przybliżenie funkcji ReLU, dzięki czemu jest różniczkowalna w~całej swojej dziedzinie, szczególnie w~punkcie zero. Dzięki temu umożliwia stabilniejsze trenowanie sieci neuronowych za pomocą metod gradientowych. Teoretycznie zakłada się lepsze zachowanie i~łatwiejsze uczenie, jednak istnieją badania, które temu zaprzeczają, sugerując użycie ReLU dla uczenia pod nadzorem.\[Softplus(a) = \ln(1 + e^{a})\]

\end{description}

Szczególnym przypadkiem są funkcje aktywacji o~zmiennym kształcie (ang. \textit{Trainable shape}). To zbiór funkcji aktywacji, których kształt jest aktualizowany podczas procesu uczenia. Pozwala to na lepsze dopasowanie funkcji aktywacji do specyficznych cech danych i~zadań, co może prowadzić do poprawy wydajności sieci neuronowej. Powstają w~wyniku dodania dodatkowych parametrów do ww. funkcji lub tworząc dodatkowe sieci neuronowe jako odrębne funkcje aktywacji.

Do odzwierciedlenia nieliniowości zbioru danych poza użyciem nieliniowych funkcji aktywacji konieczne jest dodanie \textbf{warstw ukrytych} do sieci neuronowej. Aktualizacja wag poszczególnych neuronów następuje w~wyniku algorytmu wstecznej propagacji opisanego dalej.
Warstwa w~sieci neuronowej to zbiór neuronów, które przetwarzają ten sam wektor sygnałów wejściowych (pochodzący z~poprzedniej warstwy), lecz z~użyciem różnych wag synaptycznych i~wartości progowych. 
Dla sieci neuronowych z jednym neuronie w~każdej warstwie obserwuje się, że dodawanie kolejnych warstw nie zwiększa zdolności aproksymacyjnych, a~w~granicznym przypadku prowadzi do degeneracji wyjścia do funkcji stałej.
Zarówno sieci płytkie (takie o~jednej warstwie ukrytej) jak i~głębokie mogą być uniwersalnymi aproksymatorami badanej funkcji, o~ile mają odpowiednią szerokość (liczba neuronów w~warstwie), w~przeciwieństwie do sieci bez warstw ukrytych lub do sieci, których szerokość wynosi 1 \cite{tutorial_nn_activation}.

Jedna ukryta warstwa o~odpowiedniej liczbie neuronów jest w~stanie przybliżyć większość funkcji, jednakże zwiększenie liczby warstw o~rozsądnej szerokości znacząco ułatwia to zadanie \cite{islp}.

Sieci o~pojedynczych warstwach mają ograniczoną moc obliczeniową. Ograniczenie to jest możliwe do obejścia dodając głębokość lub szerokość.
Dodawanie warstw ukrytych wiąże się z~trudnością analizy teoretycznej modelu oraz wizualizacji procesu uczenia.
Istnieją zasadniczo trzy rodzaje warstw:
\begin{enumerate}
    \item Warstwa wejściowa, której zadaniem jest przekazanie wartości argumentów aproksymowanej funkcji do warstwy ukrytej. Nie wykonują obliczeń nieliniowych, a ich wagi nie są aktualizowane w procesie uczenia.
    \item Warstwa ukryta, która nazwę swoją zawdzięcza faktowi, że obliczenia nie są widoczne z~perspektywy wejścia czy wyjścia. 
    \item Warstwa wyjściowa, ostatnia warstwa sieci, która przyjmuje jako argument wartość sygnału od ostatniej warstwy ukrytej i~przekształca na wynik końcowy sieci.
\end{enumerate}


\begin{tcolorbox}[title=Algorytm wstecznej propagacji błędu (ang. \textit{Backpropagation}),
    colback=blue!5!white,colframe=blue!75!black,
    breakable]
    
Przy założeniu, że funkcja aktywacji jest sigmoidalna, algorytm wstecznej propagacji błędu (ang. \textit{Backpropagation}) można opisać następującymi krokami:\cite{haykin2009neural}
    \begin{enumerate}
        \item \textbf{Inicjalizacja} 
        Przy założeniu, że nie ma żadnych informacji dostępnych zakłada się wagi synaptyczne oraz wartości graniczne z~rozkładu jednostajnego, którego średnia równa się zero. Wariancję rozkładu dobiera się w~taki sposób by odchylenie standardowe lokalnego indukowanego pola znajdowało się pomiędzy częścią liniową a~wysyconą funkcji sigmoidalnej.
        \item \textbf{Przekazanie danych wejściowych} 
        W~epoce $n$ aplikuje się przykład treningowy $\mathbf{x(n)}$ na warstwę wejściową, a~wartość oczekiwana $\mathbf{d(n)}$ jest przechowana do późniejszego porównania z~warstwą wyjściową.
        \item \textbf{Obliczenia „w przód” (ang. \textit{forward computation})} 
        Obliczone zostają indukowane lokalne pola dla każdej z~warstw podając wyjście jednej warstwy jako wejście kolejnej. Wartość indukowanego pola dla neuronu j-tego z~warstwy l-tej wyznacza się ze wzoru: 
        \[v_j^l(n)=\sum_i{w_{ji}^l(n)y_i^{l-1}(n)}\]
        gdzie $w_{ji}^l(n)$ to waga synaptyczna łącząca neuron i-ty z~warstwy (l-1)-tej z~neuronem j-tym z~warstwy l-tej, a~$y_i^{l-1}(n)$ to wyjście neuronu i-tego z~warstwy (l-1)-tej.
        Celem tego kroku jest wyznaczenie błędu:
        \[e_j(n) = d_j(n) - y_j^L(n)\]
        gdzie $y_j^L(n)$ to wyjście j-tego neuronu z~warstwy wyjściowej L-tej.
        \item \textbf{Obliczenia wstecz (ang. \textit{backward computation})}
        Obliczone zostają lokalne gradienty: 
        \[
            \delta_j^l(n) =
            \begin{cases}
            e_j^L(n)\,\phi'_j(v_j^l(n)) & \text{ dla neuronu } j~\text{ w~warstwie wyjściowej } L, \\
            \phi'_j(v_j^l(n))\sum_k {\delta_k^{l+1}(n)w_{kj}^{l+1}(n)} & \text{ dla neuronu } j~\text{ w~warstwie ukrytej } l,
            \end{cases}
        \]
        Gdzie $\phi'_j(\cdot)$ to pochodna funkcji aktywacji neuronu j-tego.
        Po wyznaczeniu lokalnych gradientów wykorzystuje się je do aktualizacji wag:
        \[w_{ji}^l(n+1)=w_{ji}^l(n) +\alpha\left[\Delta w_{ji}^l(n-1) \right] +\eta \delta_j^l(n)y_i^{l-1}(n) \]
        Gdzie $\eta$ to współczynnik uczenia, $\alpha$~to współczynnik momentu, a~$\Delta w_{ji}^l(n) = \eta\delta_j(n)y_i(n)$ to zmiana wagi.
        \item \textbf{Iteracja}
        Należy powtarzać kroki 3 oraz 4 wprowadzając kolejne przykłady, dopóki kryterium zatrzymania nie zostanie osiągnięte
    \end{enumerate}

\end{tcolorbox}

\subsection{Autoenkodery}

Szczególnym przypadkiem sieci neuronowych są autoenkodery. Ich głównym zadaniem jest kopiowanie danych z~wejścia na wyjście. Autoenkodery znajdują także zastosowanie w~odszumianiu danych oraz w~wydobywaniu istotnych cech danych wejściowych. Autoenkodery to kombinacja funkcji enkodera $\mathbf{h} = f(\mathbf{x})$ i~dekodera $\mathbf{r} = g(\mathbf{h})$. Enkoder wykonuje transformację wartości wejściowej na nowy format w~tzw. przestrzeni latentnej (ukrytej), a~dekoder przyjmuje nowy format jako swoje wejście i~przekształca na format pierwotny.
Trywialne rozwiązanie $g(f(\mathbf{x}))=\mathbf{x}$ nie ujawnia istotnej struktury danych. 
Autoenkodery typu \textit{undercomplete} to takie, które zmniejszają wymiar danych po przejściu przez enkoder. Proces uczenia takich autoenkoderów polega na zadaniu optymalizacyjnym funkcji straty $L(\mathbf{x},g(f(\mathbf{x})))$, która penalizuje różnicę między wejściem enkodera a~wyjściem dekodera. Autoenkoder z~liniowym dekoderze wykorzystujący błąd średniokwadratowy jako funkcję straty wykonuje te same zadania co PCA.
Autorzy \cite{deeplearningbook} wskazują jednak, że jeżeli model ma zbyt dużą „zdolność” (ang. \textit{capacity}) (np. w~wyniku wykorzystania nieliniowych enkoderów, czy zbyt skomplikowanej architektury sieci), to istnieje ryzyko, że model będzie jedynie kopiował dane bez uczenia się cech danych. Podobny problem występuje w~modelach, w~których enkoder mapuje dane wejściowe do wymiaru równego wymiarowi wejścia lub większego. Ryzyko to można zmniejszyć, utrzymując nierozbudowaną architekturę sieci lub wprowadzając specjalną funkcję straty, która karze za trywialne kopiowanie. 
Przykładem jest sparse autoenkoder będący autoenkoderem regularyzowanym. 
Różni się od autoenkodera typu undecomplete dodatkowym składnikiem funkcji straty, która oprócz kary za różnicę między wejściem a~wyjściem penalizuje również za zbyt wiele aktywnych neuronów, innymi słowy, wprowadza karę za nadmierną aktywację neuronów w~warstwie ukrytej, wymuszając rzadkość reprezentacji.

Głębokie autoenkodery o~strukturze hierarchicznie malejących warstw ukrytych (tzw. stacked autoencoders) wykazują się mniejszym błędem rekonstrukcji w~porównaniu do klasycznej analizy składowych głównych (PCA) dla tej samej liczby wymiarów przestrzeni latentnej. Wskazano również na tendencję do tworzenia wyraźnie odseparowanych klastrów, co ułatwia jakościową interpretację w~diagnostyce technicznej.

Poza optymalizacją masy komponentów elektronicznych na pokładzie samolotów systemy pokładowe wykorzystujące dane o~mniejszym wymiarze wymagają mniejszej ilości pamięci operacyjnej co skutkuje krótszym czasem wykonania (ang. \textit{runtime efficiency}) zadań i~obliczeń. To z~kolei może umożliwić wykorzystanie nowoczesnych i~tym samym bardziej skomplikowanych systemów pokładowych oferujących większe bezpieczeństwo operacyjne i~komfort użytkowania \cite{deeplearningbook}.

\subsection{Uczenie głębokie w~kontekście danych silnikowych}

Artykuł \cite{rnn_rul} opisuje predykcję żywotności silników turbowentylatorowych poprzez wykorzystanie własnościowego oprogramowania trenującego opartego na wielowarstwowym perceptronie (Multi-Layer Perceptron Neural Network). Zbiór danych jest ten sam co opisany w~rozdziale~\ref{sec:opis_danych}. Autor wykazał sensowność badania poprzez klasyfikację silników na zdrowe/zużyte przy założeniu, że pierwsze 30 cykli każdego silnika przedstawia dane typowe dla silników niewyeksploatowanych. Wspomniane oprogramowanie wykorzystuje Rozszerzony Filtr Kalmana (ang. \textit{Extended Kalman Filter Training Method}). Porównano dwie architektury: trójwarstwową i~dwuwarstwową - pierwsza wykazała większą prędkość uczenia i~dokładniejsze wyniki. Trening został przeprowadzony na systemie operacyjnym Windows XP, z~procesorem Intel o~taktowaniu 2~GHz i~osiągnięto 99,1\% dokładności klasyfikacji.

Autorzy \cite{framework_ml_health} opracowali strukturę prognozowania stanu dla rodziny silników turbowentylatorowych. Zbiór danych pochodzi z~rzeczywistych pomiarów istniejących maszyn. Zaproponowane zostały trzy rodzaje treningu w~odróżnieniu od klasycznego podejścia podziału na zbiór treningowy i~testowy. Wykorzystany model estymuje wartość temperatury gazu za turbiną podczas startu samolotu opierając się na założeniu, że jest to najlepszy wskaźnik stanu silnika z~racji na największe zapotrzebowanie na moc, duże wytężenie termiczne i~ciśnieniowe modułu turbiny. Podział został wykonany w~oparciu o~informacje dotyczące wystąpienia wizyt w~warsztatach. Autorzy zauważyli, że naprawa silników wiąże się z~poprawą wartości predykatorów żywotności w~związku z~tym trzy rodzaje podziału zostały zaproponowane. Zbiór danych jest dzielony na N+1 segmentów, gdzie N~to liczba wizyt w~warsztatach.
\begin{table}[H]
\centering
\caption{\textbf{Metody podziału zbioru danych ze względu na wizyty w~warsztacie}}
\label{tab:metody_podzialu}
\begin{tabular}{p{2.5cm} p{3cm} p{3cm} p{7cm}}
\hline
\textbf{Metoda} & \textbf{Zbiór danych treningowych} & \textbf{Zbiór danych testowych} & \textbf{Hipoteza} \\
\hline
I (one-to-many) 
& Tylko segment 1 
& Tylko segment 2 
& Segment 1 odpowiada sytuacji w~której nowy silnik zużywa się. Jest to najbardziej reprezentacyjna próba degradacji silnika. \\
\hline
II (many-to-many) 
& segment [1,N] 
& segment N+1
& Obraz pełnej sekwencji pogarszania się stanu dostarczy najwięcej informacji na ten temat zwiększając dokładność predykcji.\\
\hline
III (one-to-one) 
& segment N-1 
& segment N
& Tylko przedostatni segment obrazuje wiernie charakterystyki deterioracji.\\
\hline
\end{tabular}
\end{table}

Dodatkowo użyto dopasowania wielomianowego (ang. \textit{polynomial fitting}) w~celu odróżnienia sygnałów od trendów. Wykorzystano pierwszy segment (tj. od nowości do pierwszego remontu) każdego silnika do treningu, a~do testów wykorzystano pozostały zbiór. Ze względu na trudność związaną z~procesowaniem długich zbiorów czasowych przez klasyczne rekurencyjne sieci neuronowe (RNN) wykorzystano sieć LSTM (ang. \textit{Long Short-Term Memory} - długa pamięć krótkotrwała). RNN posiada architekturę porównywalną do łańcucha, LSTM natomiast posiada cztery warstwy, które decydują o~tym czy informacja zostanie zapomniana, przechowana czy podana dalej. 
Do oceny predykcji modeli wykorzystano błąd średniokwadratowy

\begin{equation}
    \mathrm{RMSE} = \sqrt{\frac{\sum_{i=1}^{N} \lvert \hat{y}_i - y_i \rvert^{2}}{N}}
    \label{eq:rmse}
\end{equation}
gdzie $\hat{y}$ to wartość przewidywana, a~$\mathrm{y}$ to wartość rzeczywista, N~to liczba niebrakujących danych. RMSE został obliczony dla każdego silnika, a~wyznaczona mediana i~średnia arytmetyczna zostały wykorzystane do porównania modeli.
We wszystkich podejściach podziału danych wygrywają modele liniowe (lr, lasso, ridge, en, huber, lassol, pa, sgd), których średni RMSE waha się między 3.02 – 3.94 w~zależności od użytej metody. Na drugim miejscu plasują się modele nieliniowe (knn, cart, et, svmr, ada, bag, rf, gbm), których średni RMSE wynosi 3.17 – 4.07, a~dla LSTM: 3.88 – 4.14. Średnie wartości RMSE zależą od użytej metody podziału na zbiór treningowy i~testowy.

Na podstawie modelu LSTM porównano również jakość predykcji dla różnych rodzajów podziałów dla trzech rodzin silników. Dwa na trzy rodziny najlepsze wyniki daje klasyczny podział Train-Test Split, natomiast dla jednej rodziny podejście nr 1 zwyciężyło.